% Methodology section: BrowseComp-Plus to ClimbMix projection.
% Citation keys resolve against references.bib in this directory.

\section{Benchmark Construction}
\label{sec:methodology}

We construct the benchmark by projecting questions from BrowseComp-Plus
(BCP)~\cite{Chen_etal_arXiv2025_BCP} onto the ClimbMix corpus~\cite{DiaoShizhe_etal_NeurIPS2025},
following the idea of cross-corpus relevance projection~\cite{Asadi_etal_SIGIR2011poster}. A question is
supported on ClimbMix when its gold answer can be reconstructed from a chain of intermediate facts (hops),
each of which is found in the corpus. Every retrieval step uses the same BM25 index of ClimbMix
(Pyserini~REST, \texttt{climbmix-400b}) against which the benchmark is later evaluated. The pipeline projects
each of the $830$ BCP test questions into a chain of grounded hops, checks answerability to obtain $326$
supported questions, validates them with an independent agent, and keeps the $58$ questions in which every
information-bearing clue of the question is represented by a hop grounded in the corpus, so the answer is
supported by retrieved documents. Each stage writes one output file per question, so runs are resumable
and every decision is inspectable.

\paragraph{Stage 0: Data preparation.}
We stream the BCP test split, decode each record, and build one task per question containing the question,
the gold answer, and the original supporting documents. The provided documents are used only as hints to the
intended reasoning steps; they are never counted as ClimbMix grounding.

\paragraph{Stage 1: Projection.}
An agent projects each question onto ClimbMix. From the question, its gold answer, and BCP's original
supporting documents (which ship with each question and serve as hints to the intended reasoning steps), it
first writes down the minimal chain of atomic facts, the hops, that lead from the question's constraints to
the answer, typically four to six of them. These hops capture the question's intended reasoning. The agent
then grounds each hop in the corpus: it issues BM25 queries against ClimbMix, reformulating and broadening the
query as needed, reads the retrieved documents, and attaches to each hop the ClimbMix documents that state its
fact. The provided documents are used only to recover the reasoning steps; grounding always comes from
ClimbMix. The output, for each question, is its hops together with the corpus documents that support them.

\paragraph{Stage 2: Answerability check.}
Given the projected hops and their grounding, the agent judges whether the corpus supports answering the
question. A question is supported when the grounded hops, together with reasoning over the retrieved
documents, uniquely determine the gold answer; partially supported when the answer is reachable but not
uniquely determined, or only some of the required facts are grounded; and unsupported when a decisive fact is
missing from, or contradicted by, the corpus. Answers may be reasoned out, for example by combining several
documents or converting units and dates, rather than appearing verbatim. This check is intentionally lenient:
a question is supported whenever some subset of its hops is enough to single out the answer, so individually
ungrounded or unnecessary hops do not change the label. Of the $830$ questions, $326$ are supported, $107$
partially supported, and $396$ unsupported (one further question is omitted on content-policy grounds). The
$326$ supported questions form the answerable set that the remaining stages refine.

\paragraph{Stage 3: Independent validation with piika.}
We cross-check the supported set with \textsc{piika}, an independent GPT-5.5 deep-research agent run over the
same ClimbMix BM25 endpoint. \textsc{piika} answers strictly from the documents in its context, with
parametric knowledge excluded, so a correct answer reflects the provided documents rather than the model's
memory. On its own retrieval, without any documents from us, \textsc{piika} answers $47.2\%$ ($154/326$) of
the supported questions correctly. For the remaining $172$ supported questions that it missed, we supplied the
ClimbMix documents our pipeline had identified as grounding the question's hops (the per-hop supporting
documents recorded during projection) and re-ran \textsc{piika} closed-context over exactly those documents.
It then answered $156$ of the $172$ correctly (only $16$ wrong), raising its coverage of the supported set to
$310/326$. Once given the evidence we surface, an independent agent recovers almost all of the answers it had
previously failed, confirming that those documents genuinely support them.

\paragraph{Stage 4: All-hops verification.}
The answerability check is lenient: a question passes if the answer can be reached from a sufficient subset of
hops. For a retrieval benchmark this is not enough, because a model could then reach the answer by inference,
and the inference capability could be indirectly affected by its parametric knowledge. We want the answer to
come only from the retrieved documents, so every clue in the question must be present in ClimbMix. We
decompose each answerable question into hops, one per distinguishing clue, and a fresh agent checks each hop
against the corpus: a hop is grounded when a specific ClimbMix document states its fact, kept as a verbatim
snippet, not just a document on the same topic and not the model's own knowledge. A question is kept only if
every one of its hops is grounded. Crucially, we never drop a hop to let a question pass. If even one clue in
the written question has no corpus evidence, the question is rejected, because otherwise the benchmark would
claim a grounding it does not have. An earlier version instead trimmed hops it judged unnecessary and kept a
question as long as the remaining hops still identified the answer; this let through questions whose written
clues were not all in the corpus, for example one that identifies a person's birthplace by a population-growth
statistic that ClimbMix does not contain. A hop may be labeled redundant when its fact is already implied by
another hop, but redundant hops are kept for analysis, not removed. Because BCP's full set of clues already
determines the answer uniquely, grounding every clue makes the answer uniquely derivable from the corpus, with
no missing evidence to fill in by inference. Grounding is judged by two independent scrutiny passes, and a hop
counts as grounded when either finds a verbatim document; $63$ of the $326$ answerable questions have every
listed hop grounded. Per-hop grounding can still miss a coverage gap: an
information-bearing clue may never have been written as a hop, so the hop list looks complete while a stated
condition goes unchecked. This happens most often with a specific number embedded in a phrase, such as an
interval (between $24$ and $25$ years later) or a count (married three times), where a hop captures the
surrounding event but drops the number. We therefore audit each question for coverage and remove any whose
clues are not all represented by a hop. This leaves the final benchmark of $58$ questions, each released with
its hops and the supporting ClimbMix documents. The result is strict: most BCP questions contain a clue that is
absent from ClimbMix or was never captured as a hop. The surviving set is the most corpus-answerable of all: on
its own retrieval \textsc{piika} answers $63.8\%$ ($37/58$), well above the $47.2\%$ on the answerable set. The
benchmark is thus grounded by construction, with every clue of every question represented by a hop backed by a
verbatim ClimbMix document, and cross-checked by an independent agent answering from the corpus.

\paragraph{Reproducibility.}
The projection, answerability, hop-grounding scrutiny, and the coverage audit are agentic and thus not
bit-for-bit reproducible, so we release the per-question outputs of the two scrutiny passes and the audit as
the canonical record and ship each stage as a script. The final set is regenerated deterministically from
those committed outputs, a hop counting as grounded when either pass found a verbatim document and a question
surviving only when every clue is represented by a hop, and we verified the rebuild reproduces the released set
exactly. The agentic stages can be re-run from scratch with only the ClimbMix BM25 endpoint, reproducing an
equivalent (though not identical) set.
